% =====================================================================================
% Chapter 2 · Point estimate versus posterior
%
% Built from notebooks/00_foundations.ipynb §7 and §7.1 — the passage every later chapter
% cross-references. NOT ONE NUMBER IN THIS FILE IS TYPED BY HAND: every figure in the prose
% is a macro from build/macros.tex, generated by cmp.report from the executed notebook
% (CMP_FAST=0, CMP_REAL=1). The forward-looking macros are the chapters' own: \nbSeven* is
% Chapter 9 (notebooks/07_geo_lift_synthetic_control.ipynb), \nbEight* is Chapter 10 (nb08),
% \nbSixNWeeks is Chapter 8 (nb06) — every forward claim is cited from the chapter that earns
% it, not restated from memory. The only literals below are the decision
% thresholds (0.9, 0.5) and the interval level (90 %), which are conventions stated in the
% text, and the prior scales of §2.4, which are definitions of the model rather than results.
% =====================================================================================
\chapter{Point estimate versus posterior}
\label{ch:pvp}

\begin{quote}\small
\emph{The same \nbZeroNCustomers{} customers, the same linear specification, the same
estimand, fitted twice. Ordinary least squares returns \eur{\nbZeroOlsEst} with a 90\,\%
confidence interval of $[\eur{\nbZeroOlsLo},\ \eur{\nbZeroOlsHi}]$; the Bayesian twin returns
a posterior mean of \eur{\nbZeroBayesMean} with a 90\,\% credible interval of
$[\eur{\nbZeroBayesLo},\ \eur{\nbZeroBayesHi}]$. The point estimates differ by
\eur{\nbZeroEstGap} --- \pc{\nbZeroEstGapPct} of the interval's width --- and the credible
interval is \pc{\nbZeroWidthRatioPct} wider than the confidence interval. Bayes did not
change the answer and bought nothing in precision, and this chapter says so plainly. Then a
threshold is set at \eur{\nbZeroContactCost} and the two arms are asked whether the campaign
pays. The posterior puts \nbZeroPPays{} of its mass above the cost. A bootstrap of the very
same estimator puts \nbZeroBootShare{} of its re-estimates above the same line. The numbers
agree to two decimal places; the objects do not agree at all, and only one of them is the
quantity every euro decision in this book consumes.}
\end{quote}

% =====================================================================================
\section{The decision that needs a probability}
\label{sec:pvp:decision}

Every euro decision in this book has one shape:
%
\begin{equation}
  \Prob\big(\text{effect} > \text{cost} \;\big|\; \text{data}\big)\ \begin{cases}
    \ge 0.9 & \Rightarrow\ \textsc{go},\\
    < 0.5   & \Rightarrow\ \textsc{no-go},\\
    \text{otherwise} & \Rightarrow\ \textsc{test further}.
  \end{cases}
  \label{eq:pvp:rule}
\end{equation}
%
The thresholds are a convention and can be argued with. The \emph{object} on the left cannot:
it is a probability attached to a statement about the effect, and it is what a person holding a
budget is actually asking for. \Cref{ch:foundations} produced four point estimates and four confidence
intervals of impeccable construction, and \textbf{not one of them can be fed to
\cref{eq:pvp:rule}}. That is not a computational shortcoming to be worked around with more
effort. It is structural, and this chapter is about exactly why.

The concrete decision here is small enough to hold in one hand. The email of \Cref{ch:foundations} raises
spend by a planted \eur{\nbZeroTrueDag} per contact. Suppose a contact costs
\eur{\nbZeroContactCost}. Does the campaign pay?

\begin{remark}[Why \eur{\nbZeroContactCost}, and not the \eur{\nbZeroCost} of \Cref{ch:foundations}]
Two different simulated businesses, so no cost carries over between them --- but the honest
reason the threshold is what it is deserves stating rather than dressing up. It sits about
\nbZeroCostInSds{} posterior standard deviations below the estimate, which is the only place a
comparison of this kind teaches anything. Put the cost at \eur{3} and every method prints
$1.00$; put it at \eur{8} and every method prints $0.00$; and the reader never has to ask why
the posterior's \nbZeroPPays{} and the bootstrap's \nbZeroBootShare{} are different animals.
The threshold is chosen to make the comparison bite.
\end{remark}

% =====================================================================================
\section{The confidence interval, exactly}
\label{sec:pvp:ci}

Treat the effect $\tau$ as a fixed, unknown constant. The \emph{data} are what is random: the
\nbZeroNCustomers{} customers in hand are one draw from a population that could have handed over
a different \nbZeroNCustomers. An interval estimator is a pair of functions
$[\,L(D),\, U(D)\,]$ of the dataset $D$, and it has 90\,\% coverage when
%
\begin{equation}
  \Prob_{\,D \,\sim\, P_\tau}\Big(\,L(D) \;\le\; \tau \;\le\; U(D)\,\Big) \;=\; 0.90
  \qquad\text{for every possible value of } \tau .
  \label{eq:pvp:coverage}
\end{equation}
%
Read where the probability sits. It sits on $D$ --- on the dataset that could have come out
otherwise. It does \textbf{not} sit on $\tau$, which is not a random variable at all in this
calculus and therefore cannot carry a probability. And note the quantifier at the end: coverage
must hold \emph{for every} $\tau$, which is what makes \cref{eq:pvp:coverage} a property the
procedure can be held to before any data arrive.

The consequence is uncomfortable and completely rigorous. Once the data are in hand and the
interval is printed, the printed interval either contains $\tau$ or it does not. There is no
``90\,\% chance'' left anywhere in it, exactly as a coin already flipped and covered by a hand
is not ``50\,\% heads'': one is simply ignorant of a settled fact, and ignorance is not
probability in this framework. \textbf{Ninety per cent describes the factory, not the particular
interval that came off the line.}

That is not a rhetorical flourish; it is the only thing \cref{eq:pvp:coverage} can be audited
against, and \Cref{ch:foundations} performed the audit in the only way available: rebuild the interval on
fresh samples and count. Across \nbZeroNSeeds{} fresh confounded panels, the AIPW 90\,\%
interval contained the planted effect on \nbZeroAipwCovered{} of them ---
\pc{\nbZeroAipwCovPct}. The factory works. What the factory cannot do, at any level of effort,
is tell you the probability that \emph{this} interval, the one on the page, contains the truth.

\begin{proposition}[What a confidence interval licenses, and what it does not]
\label{prop:pvp:ci}
From \cref{eq:pvp:coverage} one may say: \emph{if this procedure were repeated on fresh samples,
90\,\% of the intervals it produced would contain the true effect.} One may not say: \emph{there
is a 90\,\% probability that the true effect lies in $[\eur{\nbZeroOlsLo},\ \eur{\nbZeroOlsHi}]$.}
The second sentence assigns a probability to $\tau$, and the framework that produced the first
sentence contains no such probability to assign. In particular, no manipulation of a confidence
interval, a $p$-value or a standard error yields $\Prob(\tau > c)$: there is no distribution over
$\tau$ whose mass above $c$ could be taken.
\end{proposition}

\Cref{prop:pvp:ci} is the whole difficulty, and it is worth being blunt about how it is handled
in practice: the second sentence is what almost everyone means when they read the first one
aloud in a meeting. The sentence people want is a Bayesian sentence. The question is whether they
are willing to pay for it.

% =====================================================================================
\section{The credible interval, exactly}
\label{sec:pvp:cri}

The price of admission is a distribution over $\tau$ itself. Assert a \emph{prior}
$p(\tau)$ --- what is believed about the effect before the data --- write a \emph{likelihood}
$p(D \mid \tau)$ --- what the model says the data would look like for each value of $\tau$ ---
and Bayes' theorem returns the \emph{posterior}
%
\begin{equation}
  p(\tau \mid D) \;=\; \frac{p(D \mid \tau)\, p(\tau)}{\int p(D \mid \tau')\, p(\tau')\, d\tau'} .
  \label{eq:pvp:bayes}
\end{equation}
%
Now $\tau$ \emph{is} a random variable, and a 90\,\% \emph{credible interval} is any interval
$[a, b]$ with
%
\begin{equation}
  \Prob\Big(\,a \le \tau \le b \;\Big|\; D,\ \text{model},\ \text{prior}\,\Big) \;=\; 0.90 .
  \label{eq:pvp:cri}
\end{equation}
%
Compare \cref{eq:pvp:cri} with \cref{eq:pvp:coverage}, because the difference is the chapter. In
\cref{eq:pvp:coverage} the randomness is in $D$ and the statement is about a procedure. In
\cref{eq:pvp:cri} the randomness is in $\tau$ and the statement is about the effect, conditional
on the data actually collected. \Cref{eq:pvp:cri} is the sentence everyone believes they are
saying. It is a legitimate sentence. It is legitimate \emph{only after} the two conditioning
terms on the right of the bar have been paid for: a model, and a prior --- both of which must
then be defended, and both of which \cref{sec:pvp:notbuy} attacks.

And \cref{eq:pvp:cri} is not the point of the exercise. The point is that the same posterior
answers a question \cref{eq:pvp:coverage} cannot even be asked:
%
\begin{equation}
  \Prob\big(\tau > c \;\big|\; D\big) \;=\; \int_{c}^{\infty} p(\tau \mid D)\, d\tau
  \;\approx\; \frac{1}{S}\sum_{s=1}^{S} \mathbf 1\big[\,\tau^{(s)} > c\,\big],
  \label{eq:pvp:pgt}
\end{equation}
%
where $\tau^{(1)}, \dots, \tau^{(S)}$ are posterior draws. \Cref{eq:pvp:pgt} is one line of
arithmetic over the draws, and it is precisely the left-hand side of the decision rule
\cref{eq:pvp:rule}. \textbf{This is the only thing the Bayesian layer buys in this chapter, and
it is enough.}

% =====================================================================================
\section{The two arms, on identical data}
\label{sec:pvp:arms}

The demonstration is deliberately anticlimactic, and it is arranged so that \emph{nothing}
differs between the arms except the regime of inference. The classical arm is rung~2 of
\Cref{ch:foundations}'s ladder: ordinary least squares of spend on the email and on loyalty, with an HC1
standard error. The Bayesian arm is the same specification with priors bolted on:
%
\begin{align}
  y_i \;\big|\; \alpha, \tau, \beta, \sigma
    &\;\sim\; \mathcal N\big(\alpha + \tau\,T_i + \beta\,L_i,\ \sigma^{2}\big),
    \label{eq:pvp:lik}\\
  \alpha &\;\sim\; \mathcal N(20,\, 10^{2}), \qquad
  \tau \;\sim\; \mathcal N(0,\, 10^{2}), \label{eq:pvp:prior_tau}\\
  \beta &\;\sim\; \mathcal N(0,\, 10^{2}), \qquad
  \sigma \;\sim\; \text{HalfNormal}(10). \label{eq:pvp:prior_sig}
\end{align}
%
\Cref{eq:pvp:lik} is the likelihood, and it is \emph{exactly} rung~2's outcome model --- same
regressors, same functional form, same conditional mean. The rest are priors, each
defensible in a clause: all are weakly informative on the euro scale of the data ($\alpha$
centred near a typical spend of about \eur{20}; $\beta$ and $\sigma$ wide enough for anything
plausible), and $\tau \sim \mathcal N(0, 10^2)$ is sceptical-but-open --- centred on ``the email
does nothing'', yet wide enough to learn a \eur{30} effect if the data insist on one. The
identification is untouched: loyalty is in \cref{eq:pvp:lik} for exactly the reason
\cref{as:fnd:backdoor} put it in rung~2, and the sampler adds no causal content whatsoever.

One asymmetry between the arms is deliberate and should not be smoothed over, because it is the
axis the rest of the book turns on. The arms share a conditional mean; they do \emph{not} share an
error structure. \Cref{eq:pvp:lik} commits to one --- errors that are Normal, independent and of
constant variance --- while the classical arm's HC1 covariance declines to commit, estimating the
error scale from the residuals instead of asserting it (the remark in \cref{sec:pvp:spine} is about
that refusal). On these data the commitment is free: the simulator's errors really are
homoskedastic, so HC1 and the classical standard error agree and nothing turns on the difference.
\Cref{sec:pvp:spine} is what happens when the commitment is not free.

Inference is by the No-U-Turn sampler \citep{hoffman2014} in PyMC \citep{salvatier2016}:
\nbZeroNDraws{} post-warmup draws across \nbZeroNChains{} chains, with $\hat R =
\nbZeroRhat$, a minimum bulk effective sample size of \nbZeroEss{} \citep{vehtari2021}, and
\nbZeroDivergences{} divergent transitions. The chains have converged; nothing below is a
sampler artefact.

\input{build/tables/nb00_side}

\Cref{tab:pvp:side} is the result, and it is a flat one. The classical point estimate is
\eur{\nbZeroOlsEst}; the posterior mean is \eur{\nbZeroBayesMean}. They differ by
\eur{\nbZeroEstGap} --- \pc{\nbZeroEstGapPct} of the confidence interval's width, which is to say
by nothing. The confidence interval is $[\eur{\nbZeroOlsLo},\ \eur{\nbZeroOlsHi}]$, width
\eur{\nbZeroOlsWidth}. The credible interval is $[\eur{\nbZeroBayesLo},\ \eur{\nbZeroBayesHi}]$,
width \eur{\nbZeroBayesWidth} --- \pc{\nbZeroWidthRatioPct} wider. The doubly-robust rung, which
uses a different estimator entirely, lands at \eur{\nbZeroAipwEst} with a bootstrap interval of
width \eur{\nbZeroAipwWidth}.

State the finding without decoration, because a book that oversells its own method is not worth
reading. \textbf{On this problem, Bayes did not change the answer, and bought nothing in
precision.} The credible interval is not tighter than the confidence interval; it is marginally
wider. Anyone who suspected the Bayesian arm of smuggling in the answer it wanted may put the
suspicion down. Both arms are reading the same data through the same specification, honestly, and
arriving at the same place.

The only column of \cref{tab:pvp:side} that differs is the last one --- the \emph{kind} of
interval --- and that column is the entire chapter.

% =====================================================================================
\section{Why they agree: Bernstein--von Mises}
\label{sec:pvp:bvm}

The agreement in \cref{tab:pvp:side} is not luck, and it is not a property of this dataset. It is
a theorem, and knowing which theorem matters, because the theorem's conditions are exactly the
places where later chapters break.

\begin{proposition}[Bernstein--von Mises, informally]\label{prop:pvp:bvm}
Let the model be correctly specified and regular, the parameter finite-dimensional, and the prior
positive and continuous in a neighbourhood of the truth. Then as $n \to \infty$ the posterior
distribution of $\sqrt{n}(\tau - \hat\tau_{\text{MLE}})$ converges to the same Gaussian
distribution as the sampling distribution of $\sqrt{n}(\hat\tau_{\text{MLE}} - \tau)$
\citep[Ch.~10]{vandervaart1998}. The prior washes out; the posterior concentrates where the
maximum-likelihood estimate does; and a $1-\alpha$ credible interval is asymptotically a
$1-\alpha$ confidence interval.
\end{proposition}

That is why \cref{tab:pvp:side} is boring, and it is why this book expects the two arms to agree
in most chapters and says so up front. With \nbZeroNCustomers{} customers, a four-parameter linear
model, weakly informative priors and a well-specified likelihood, \cref{prop:pvp:bvm} is not an
asymptotic aspiration --- it has already arrived, to two decimal places.

Three conditions in \cref{prop:pvp:bvm} are load-bearing, and it is worth naming them now because
the rest of the book is, in a sense, a tour of their failures.

\begin{itemize}[leftmargin=*]
  \item \textbf{``Correctly specified''.} If the likelihood asserts an error structure the data do
        not have, the posterior still concentrates --- but around the wrong \emph{scale}. The
        misspecified analogue of \cref{prop:pvp:bvm} \citep{kleijn2012} converges to a Gaussian
        whose variance is the sandwich (Huber--White) covariance, \emph{not} the inverse Fisher
        information the posterior actually reports. \Cref{sec:pvp:spine} is that failure, measured
        in euros.
  \item \textbf{``Enough data''.} Thin data is where the prior stops washing out. That is not a
        defect; it is the point. \Cref{ch:mmm}'s media-mix model fits saturation curves to
        \nbSixNWeeks{} weeks, and there the priors carry real weight and must be defended rather
        than defaulted.
  \item \textbf{``Regular''.} Hierarchical shrinkage across many thin segments,
        boundary parameters, and nonlinear functionals of the parameters all sit outside the
        theorem's comfort. They are also, not coincidentally, where the Bayesian machinery earns
        the most.
\end{itemize}

% =====================================================================================
\section{The bootstrap look-alike}
\label{sec:pvp:bootstrap}

Now the decision. At a cost of \eur{\nbZeroContactCost}, \cref{eq:pvp:pgt} over the posterior
draws gives
%
\begin{equation}
  \Prob\big(\tau > \eur{\nbZeroContactCost} \,\big|\, D\big) \;=\; \nbZeroPPays ,
  \label{eq:pvp:ppays}
\end{equation}
%
and the rule \cref{eq:pvp:rule} returns \textbf{\nbZeroVerdict}. That is a probability about the
effect, computed from a distribution over the effect, and it is the number the business asked for.

The classical arm can manufacture something that looks exactly like it, and it is important to
build the look-alike rather than dismiss it, because in practice it is what people do. Bootstrap
the OLS estimator: draw \nbZeroNBoot{} resamples of the \nbZeroNCustomers{} customers with
replacement, refit rung~2 on each, and count the share of re-estimates above the cost line. The
answer is \nbZeroBootShare.

The same number, to two decimal places. The two distributions are drawn on top of each other in
\cref{fig:pvp:overlay}, and they are very nearly the same picture.

\input{build/floats/nb00_point_vs_post}

\subsection*{Why they coincide}

They coincide for the reason \cref{sec:pvp:bvm} gave. The bootstrap distribution of $\hat\tau$
approximates the \emph{sampling} distribution of $\hat\tau$ \citep{efron1979}; the posterior
approximates the \emph{Bayesian} distribution of $\tau$; and \cref{prop:pvp:bvm} says that in the
friendly regime --- large $n$, weak priors, a symmetric and well-specified likelihood --- these
two are asymptotically the same Gaussian, merely reflected about the estimate. A symmetric
Gaussian reflected about its centre is itself. So the two integrals over the region $\{\,\cdot >
c\,\}$ agree, here to within \nbZeroBootGap.

\subsection*{Why coincidence is not identity}

And yet they are different objects, and the difference is not philosophical bookkeeping.

The bootstrap distribution answers: \emph{if I re-ran the sampling of customers many times, how
would my estimate $\hat\tau$ scatter?} Its support is a set of possible \emph{estimates}. It is
constructed entirely from datasets that were never collected, and it contains no statement about
$\tau$ --- $\tau$ is a fixed constant throughout the entire construction, exactly as in
\cref{eq:pvp:coverage}. To say ``\nbZeroBootShare{} of the bootstrap re-estimates exceed the cost,
therefore there is a \nbZeroBootShare{} probability that the campaign pays'' is to change the
subject mid-sentence, from a statement about $\hat\tau$ to a statement about $\tau$, without ever
having introduced the object that would license it.

The posterior answers: \emph{given the data I actually have, this model and this prior, how
plausible is each value of $\tau$?} Its support is a set of possible \emph{effects}. Its mass above
$c$ is $\Prob(\tau > c)$ by construction, not by resemblance.

\Cref{tab:pvp:probability} lays the three objects of this chapter side by side, and the middle
column --- \emph{what it is a probability of} --- is the only column that matters.

\input{build/tables/nb00_probability}

\subsection*{Where they separate}

The near-coincidence in \cref{fig:pvp:overlay} is precisely why the confusion is so durable and so
hard to kill: for years, in the friendly case, reading a bootstrap as a belief costs nothing at
all. It stops being free exactly where the money is.

\begin{itemize}[leftmargin=*]
  \item \textbf{Thin data.} When $n$ is small, the prior does not wash out and there is no theorem
        left. The bootstrap of a thin sample resamples its own idiosyncrasies; the posterior
        regularizes toward a stated prior. They do not agree, and it is no longer obvious which
        one you would rather defend --- but only one of them has told you what it assumed.
  \item \textbf{Hierarchical shrinkage.} Partial pooling across segments makes each segment's
        estimate depend on the others through a hyperparameter --- as \Cref{ch:pe}'s hierarchical
        elasticities and \Cref{ch:mmm}'s media-mix model both do. Shrinkage itself is \emph{not} the
        dividing line, and \Cref{ch:seg} is careful to say so: it pools its segments with the
        empirical-Bayes (James--Stein) estimator, a closed form with no prior and no sampler. The
        dividing line is what comes out. Even after shrinking, the classical arm holds a point
        estimate and a standard error for segment $s$; it holds no distribution over $\tau_s$, so it
        cannot return $\Prob(\tau_s > \text{cost})$ --- and a bootstrap of segment $s$ alone cannot
        manufacture one, because it resamples that segment's own idiosyncrasies, which is precisely
        what the pooling was there to discount.
  \item \textbf{Nonlinear functionals.} Return on ad spend, a saturation curve's headroom, the
        value of a targeting policy: each is a nonlinear function $g(\theta)$ of the parameters.
        Uncertainty must be \emph{propagated} through $g$, not reported alongside it, and the
        posterior does this by construction --- push every draw through $g$ and read the resulting
        distribution. The classical arm needs a delta method (which linearizes, and is wrong
        exactly where $g$ curves) or a bootstrap of the whole pipeline (which is often what the
        Bayesian sampler was doing anyway, but slower and without the prior stated). \Cref{ch:mmm} does
        this for ROAS; \Cref{ch:up} does it for a policy, in euros.
  \item \textbf{Any decision at all.} \Cref{eq:pvp:rule} is a functional of the posterior. There is
        no classical object to feed it.
\end{itemize}

\begin{remark}
There is an honest escape hatch, and it should be named. Under a flat prior and a regular
likelihood, the posterior and the bootstrap coincide \emph{by theorem}, and one may then read the
bootstrap as a belief with a clear conscience. But notice what has happened: a prior was asserted,
a likelihood was asserted, and the resulting object was interpreted as a distribution over the
parameter. That is Bayesian inference. The escape hatch is real, and it leads into the building.
\end{remark}

% =====================================================================================
\section{What the posterior does not buy}
\label{sec:pvp:notbuy}

Three things, and this section exists because a chapter that only listed the advantages would be
an advertisement.

\subsection*{It does not buy identification}

This is the important one, and it is the same lesson \Cref{ch:foundations} taught with four frequentist rungs,
now taught again with a sampler --- which is the point. Take the model of \cref{eq:pvp:lik}, keep
the priors, keep the sampler, and drop \emph{one covariate}: loyalty. The backdoor path of
\Cref{ch:foundations} re-opens, and the posterior does what a posterior does. It concentrates.

It concentrates on \eur{\nbZeroNoadjMean} --- a jump of \eur{\nbZeroNoadjJump}, landing exactly on
the confounded naive answer --- against a truth of \eur{\nbZeroTrueDag}. And it does so inside an
interval only \pc{\nbZeroNoadjWiderPct} wider than the correct model's, which is a gap no reader on
earth would take as a warning. Its $\hat R$ is fine. Its effective sample size is fine. It has no
divergences. It is, as a piece of Bayesian inference, impeccable. It is simply answering a
different question from the one that was asked.

\input{build/floats/nb00_confident_wrong}

\Cref{fig:pvp:wrong} is the most important figure in this book, and the sentence it earns is:
\textbf{being confidently wrong is not punished with a wide interval}. A posterior quantifies
uncertainty about the parameters of the model it is handed. It has no opinion whatever about
whether that model identifies a causal effect. The prior did not fail here --- \emph{the graph
did}, and there is no prior on $\tau$ that encodes ``you forgot a confounder''.

This is the discipline the whole book runs on, and it is worth one line of its own: \textbf{Bayes
and machine learning sharpen estimation; they never touch identification.} The DAG decides which
variables enter the model. The likelihood and the prior only quantify uncertainty about their
coefficients. Every chapter that follows settles identification --- assumptions named, numbered,
and tested --- \emph{before} it admires a posterior.

\subsection*{It does not buy precision}

Already established, and worth repeating because it is so often assumed otherwise. The credible
interval in \cref{tab:pvp:side} is \pc{\nbZeroWidthRatioPct} \emph{wider} than the HC1 confidence
interval on identical data. When both arms are reading the same likelihood, there is no free
information for either of them to find.

\subsection*{The prior does almost nothing}

The last of the three, and the one that most surprises the sceptic in the room. Sweep the prior
standard deviation on $\tau$ across a factor of \nbZeroPriorRatio{} --- from
$\tau \sim \mathcal N(0, 2^2)$, a sceptic who all but rules out an effect of the planted size,
through the default $\mathcal N(0, 10^2)$, to $\mathcal N(0, 50^2)$, which is flat across any
plausible range. The posterior mean never moves more than \eur{\nbZeroPriorShift} from the default
fit's \eur{\nbZeroBayesMean}, spanning \eur{\nbZeroPriorTightMean} under the sceptic to
\eur{\nbZeroPriorFlatMean} under the flat prior (\cref{tab:pvp:prior}, \cref{fig:pvp:prior}). The
posterior standard deviation does not move at all to two decimal places.

\input{build/tables/nb00_prior}
\input{build/floats/nb00_prior_sweep}

\textbf{The likelihood is the whole ballgame.} This is the reply to ``your priors are driving the
answer'', and it has the virtue of being a measurement rather than an assurance. Which is exactly
why the same sweep is the instrument that would reveal the \emph{opposite} verdict on thin data,
where the prior does carry weight and must therefore be argued for --- \Cref{ch:mmm}'s media-mix
model, on \nbSixNWeeks{} weeks and collinear channels, is where that bill comes due.

And it sets up the finding the rest of the book is built around. If the prior does almost nothing
and the likelihood does everything, then \textbf{when a Bayesian analysis goes wrong, look at the
likelihood}.

% =====================================================================================
\section{The book's spine}
\label{sec:pvp:spine}

Four claims, stated here and earned across the twelve chapters that follow. Each is a claim that
can be disagreed with, and each is cited forward to the chapter that pays for it.

\subsection*{1. Classical and Bayesian usually agree}

This chapter is the first instance and the cleanest: \eur{\nbZeroEstGap} apart, on a
\eur{\nbZeroTrueDag} effect. \Cref{ch:geo}, on an entirely different estimator (synthetic control on a
geo panel), finds its two arms \pc{\nbSevenEstDiffPct} of the effect apart. \Cref{prop:pvp:bvm} is
why, and the honest expectation to carry into every chapter is that the sampler will confirm the
point estimate rather than move it. When it does not, something has been learned --- and what has
been learned is almost never ``the paradigm was wrong''.

\subsection*{2. Where Bayes lost, the likelihood was misspecified --- not the paradigm}

This is the claim the book most needs to survive contact with a hostile reader, and it is settled
by a pair of chapters that fail in \emph{opposite directions}.

\Cref{ch:geo} estimates the total lift of a television campaign over \nbSevenNPost{} weeks. Its
Bayesian model, as first written, asserts an \textbf{iid} likelihood over the weekly errors ---
the default, and the one nearly everyone writes. Under independence, the
standard deviation of an $n$-week sum scales as $\sqrt{n}\,\sigma$; under correlation, the variance
of the sum collects every covariance term,
%
\begin{equation}
  \Var\Big(\sum_{t=1}^{n} e_t\Big) \;=\; \sigma^2\Big(n + 2\sum_{k=1}^{n-1}(n - k)\,\rho_k\Big)
  \;\xrightarrow[\ \rho_k \to 1\ ]{}\; n^2\sigma^2 \;\gg\; n\,\sigma^2 .
  \label{eq:pvp:n2}
\end{equation}
%
The errors there have a measured lag-one autocorrelation of \nbSevenRhoOne{} --- they are a
persistent offset, not fresh weekly noise --- and the likelihood set every $\rho_k$ in
\cref{eq:pvp:n2} to zero. The consequence is exactly what \cref{eq:pvp:n2} predicts: fitted with its
first, iid likelihood, the model prices the standard deviation of the \nbSevenNPost-week total at
\keur{\nbSevenIidSd} when the true sampling variability of that total, measured across fresh panels,
is \keur{\nbSevenSeedSd}. It underprices the uncertainty by a factor of \nbSevenUnderprice, and its
nominal 90\,\% interval covers the truth on \pc{\nbSevenTotalCov} of panels, against
\pc{\nbSevenRiCov} for the design-based randomization interval on the same panels.

That, alone, could be read as an indictment of Bayesian inference. It is not, and there are two
proofs. The first is that \Cref{ch:did}'s difference-in-differences errs in the \textbf{opposite}
direction: its pooled posterior interval comes back \nbEightBayesOverCluster$\times$ \emph{wider}
than the clustered classical one --- a half-width of \nbEightPooledOverScatter$\times$ the
estimator's true seed-to-seed scatter, where a calibrated 90\,\% interval sits at
\nbEightCalibratedRatio$\times$ --- because the store baseline was left in $\sigma$. Too wide, not
too narrow. A paradigm cannot be systematically overconfident and systematically underconfident at
the same time. A \emph{specification} can be wrong in either direction, and specifications are what
differ between those two chapters.

The second proof is the repair, and \Cref{ch:geo} performs it. The diagnosis names a defect of the
likelihood, so the correction is applied to the likelihood: give the model the error structure the
data actually have --- an AR(1) residual, as in Bayesian structural time series
\citep{brodersen2015}. Nothing else moves. Same estimand, same donor pool, same priors, same
sampler. The model recovers a posterior lag-one correlation of \nbSevenRhoPost{}, and its nominal
90\,\% interval on the \nbSevenNPost-week total goes from covering at \pc{\nbSevenTotalCov} to
covering at \pc{\nbSevenTotalCovArone} --- level with the \pc{\nbSevenRiCov} of the design-based
randomization interval, which assumed nothing about the errors in the first place. The point
estimate barely moves at all: the two arms of \Cref{ch:geo} finish \pc{\nbSevenEstDiffPct} of the effect
apart. \textbf{Repair the likelihood and the frequentist coverage comes back} --- which is the
falsifiable prediction ``specification, not paradigm'' makes, tested and confirmed in euros. The
alternative repair is to correct the posterior for misspecification directly: \citet{mueller2013}
shows that a posterior computed under a misspecified likelihood can be rescaled by a sandwich
covariance so that its intervals recover their frequentist coverage, which is the right instrument
when a posterior's \emph{centre} is fine and only its \emph{scale} is wrong.

\begin{remark}[Why the classical toolkit wins those fights]
Clustered and HAC standard errors \citep{newey1987, bertrand2004}, the bootstrap, and randomization
inference all have one property in common: they decline to assume a shape for the errors, and go and
\emph{measure} one instead. That agnosticism is the whole reason they win when a model is wrong. It
is not that frequentism is deeper. It is that it committed to less.
\end{remark}

\subsection*{3. The prior does almost nothing; the likelihood is the whole ballgame}

Measured in \cref{tab:pvp:prior}: a twenty-five-fold prior sweep moves the answer by
\eur{\nbZeroPriorShift}. The corollary is the debugging rule of this book. When a posterior
disagrees with a design-based interval, do not adjust the prior. Interrogate the likelihood, and in
particular interrogate what it assumes about the \emph{correlation structure of the errors} --- the
one assumption that is invisible in a well-fitting model and fatal in a sum.

\subsection*{4. What survives everywhere is $\Prob(\text{effect} > \text{cost})$}

The classical arm cannot produce it (\cref{prop:pvp:ci}); every euro decision in this book consumes
it (\cref{eq:pvp:rule}). That asymmetry is why the Bayesian arm is in the book at all.

But the claim must be made with its own qualification attached, because \Cref{ch:geo} supplies the
uncomfortable case: \textbf{the decision quantity inherits the scale of whatever interval it is read
from}. Its television campaign is decided on the repaired AR(1) posterior --- the one that covers at
\pc{\nbSevenTotalCovArone} --- which puts \nbSevenPGo{} of its mass above the cost of the campaign
and so returns \textbf{\nbSevenDecision} under \cref{eq:pvp:rule}. Re-scale that posterior to the
design-based standard deviation, wider than its own, and the probability moves to
\nbSevenPGoRescaled{} --- it \emph{rises}, because the posterior mean sits \emph{below} the cost
line, so widening the distribution moves mass across it from the wrong side. The call survives
(\nbSevenDecisionRescaled), and it survives because the likelihood was repaired first, not because
the arithmetic was kind. The general lesson is direction-free and worth stating in that form: a
mispriced scale does not reliably flatter the campaign, it distorts $\Prob(\tau > \text{cost})$
toward whichever verdict the mispricing happens to favour --- and nothing inside the posterior
announces which. Only the design-based apparatus does.

Which is why the guidance the book converges on is a division of labour rather than a choice of
camp.

\begin{quote}
\textbf{Use the posterior for the decision, and design-based inference for your confidence in it.}
\end{quote}

The posterior is the only object that answers the question the business asked. The design-based
apparatus --- randomization inference, placebo distributions, cluster-robust and HAC standard
errors, a bootstrap that assumes nothing about the errors --- is the only apparatus that can tell
you whether the posterior's width deserves to be believed. Compute both. Where they agree, ship.
Where they disagree, the likelihood is lying to you, and the design-based number is the one to
trust.

% =====================================================================================
\section{Takeaways}
\label{sec:pvp:take}

\begin{enumerate}[leftmargin=*, itemsep=.45em]
  \item \textbf{A confidence interval is a property of the procedure.} \Cref{eq:pvp:coverage} puts
        the probability on the \emph{data}, not on the effect. The printed interval either contains
        the truth or it does not, and ninety per cent describes the factory rather than the
        interval that came off the line. \Cref{ch:foundations} audited that factory in the only way it can be
        audited: \nbZeroAipwCovered{} of \nbZeroNSeeds{} fresh samples, \pc{\nbZeroAipwCovPct}.

  \item \textbf{A credible interval is a probability statement about the effect} --- and it is the
        sentence everyone believes they are saying when they read a confidence interval aloud
        (\cref{eq:pvp:cri}). It is legitimate. Its price of admission is a model and a prior, both
        of which must then be defended.

  \item \textbf{On easy data the two arms agree, and this chapter says so plainly.} Point estimates
        \eur{\nbZeroEstGap} apart; a credible interval \pc{\nbZeroWidthRatioPct} \emph{wider} than the
        confidence interval. \textbf{Bayes did not change the answer and bought nothing in
        precision.} That is Bernstein--von Mises (\cref{prop:pvp:bvm}), not luck, and it is the
        expected result in most chapters of this book.

  \item \textbf{The bootstrap look-alike is a different animal wearing the same clothes.}
        \nbZeroPPays{} of the posterior sits above the cost; \nbZeroBootShare{} of the bootstrap
        re-estimates sit above the same line. The first is a probability about the \emph{effect};
        the second is the scatter of the \emph{estimator} across datasets that were never collected
        (\cref{tab:pvp:probability}). They agree in the friendly case --- which is exactly why the
        error survives in industry --- and they separate on thin data, under hierarchical shrinkage,
        and through any nonlinear functional such as ROAS or the value of a policy.

  \item \textbf{The posterior does not buy identification.} Drop one covariate and the posterior
        concentrates, serenely and with excellent diagnostics, on \eur{\nbZeroNoadjMean} against a
        truth of \eur{\nbZeroTrueDag} --- in an interval only \pc{\nbZeroNoadjWiderPct} wider than
        the correct model's (\cref{fig:pvp:wrong}). Being confidently wrong is not punished with a
        wide interval. The graph decides what enters the model; Bayes only prices the coefficients.

  \item \textbf{The prior does almost nothing; the likelihood is the whole ballgame.} A
        twenty-five-fold prior sweep moves the posterior mean by \eur{\nbZeroPriorShift}
        (\cref{tab:pvp:prior}). So when a Bayesian analysis goes wrong, the suspect is the
        likelihood --- and above all what it assumed about the correlation of the errors
        (\cref{eq:pvp:n2}), which is invisible in a fit and fatal in a sum.

  \item \textbf{Where Bayes lost, the likelihood was misspecified --- not the paradigm.} \Cref{ch:geo}'s
        iid likelihood underprices a \nbSevenNPost-week sum by \nbSevenUnderprice$\times$ and covers
        at \pc{\nbSevenTotalCov}; \Cref{ch:did}'s errs the \emph{other} way, returning an interval
        \nbEightBayesOverCluster$\times$ too wide. A paradigm
        cannot be overconfident and underconfident at once; a specification can. And the diagnosis
        is testable: put an AR(1) residual into \Cref{ch:geo}'s likelihood and nothing else, and its
        coverage returns to \pc{\nbSevenTotalCovArone} --- level with the design-based
        \pc{\nbSevenRiCov} --- while the point estimate moves by \pc{\nbSevenEstDiffPct} of the
        effect. Repair the likelihood \citep{brodersen2015}, or repair the posterior
        \citep{mueller2013}. Do not repair your faith in Bayes.

  \item \textbf{The rule.} \textbf{Use the posterior for the decision, and design-based inference
        for your confidence in it} --- and remember that $\Prob(\text{effect} > \text{cost})$
        inherits the \emph{scale} of the interval it was read from, in whichever direction that
        scale is wrong.
\end{enumerate}

\notebooknote{%
  This chapter is \S7 and \S7.1 of \code{notebooks/00\_foundations.ipynb}; \Cref{ch:foundations} is the rest of
  it. The data are \code{cmp.dgp.dag\_control\_demo} (\nbZeroNCustomers{} customers, planted effect
  \eur{\nbZeroTrueDag}); the classical arm is \code{cmp.classical.ols} with an HC1 covariance, the
  Bayesian arm is the PyMC model of \crefrange{eq:pvp:lik}{eq:pvp:prior_sig}, and the seeds are fixed
  in the notebook. Every \code{nbZero} number here comes from the full run (\code{CMP\_FAST=0}); every
  \code{nbSeven} number in \cref{sec:pvp:spine} is \Cref{ch:geo}'s, emitted by
  \code{notebooks/07\_geo\_lift\_synthetic\_control.ipynb} --- the forward claims of this chapter are
  injected from the chapter that earns them, not restated from memory.}
